\subsection{Exact waves over a linear background}\label{cc:sec:affine}
We first work on $\R^3$, where linear velocity fields give local models for the
periodic construction. Let $D(t)$ be smooth and trace free, with
$D'+D^2$ symmetric. Then $Dx$ solves the Euler equations with pressure
$-x\cdot(D'+D^2)x/2$. Choose a smooth odd $1$-periodic profile
$F$ with $F'(0)=1$, and let $P'=F$, $\av P=0$.
Here $\av H=\int_0^1H(s)\,\mathrm ds$ is the phase average.

\begin{lemma}\label{cc:lem:wave}
Let the wave vector $k(t)\ne0$ and the amplitude vector $\mathbf a(t)$,
functions of time alone, satisfy
\begin{equation}\label{cc:eq:wave-ode}
 k'=-D^Tk,\qquad
 \mathbf a'=-D\mathbf a+p_{\rm w} k,\qquad
 p_{\rm w}=\frac{2k\cdot D\mathbf a}{|k|^2},\qquad
 k\cdot\mathbf a=0.
\end{equation}
Then $Dx+\mathbf aF(k\cdot x)$ is divergence free and solves the Euler equations on $\R^3$ with
pressure $-x\cdot(D'+D^2)x/2-p_{\rm w} P(k\cdot x)$.
If $k(t_0)\ne0$, $\mathbf a(t_0)\ne0$ and
$k(t_0)\cdot\mathbf a(t_0)=0$, then both vectors remain nonzero and
perpendicular for every $t\ge t_0$ in the interval of existence of $D$.
\end{lemma}
\begin{proof}
The equation $k'=-D^Tk$ transports the phase by $Dx$, and
perpendicularity cancels the self-advection:
$(\mathbf aF(k\cdot x)\cdot\nabla)\,\mathbf aF(k\cdot x)
=(k\cdot\mathbf a)FF'\mathbf a=0$. The remaining
acceleration $(\mathbf a'+D\mathbf a)F=p_{\rm w} kF$ is the gradient of
$p_{\rm w} P(k\cdot x)$.
The choice of $p_{\rm w}$ gives
$(k\cdot\mathbf a)'=-2k\cdot D\mathbf a+p_{\rm w}|k|^2=0$.
For a given $D$, the equation for $k$ is linear and has an invertible
fundamental matrix. Hence $k(t_0)\ne0$ implies $k(t)\ne0$. With this
$k$, the equation
$\mathbf a'=-(\Id-2\frac{k}{|k|}\frac{k^T}{|k|})D\mathbf a$
is also linear with an invertible fundamental matrix, so
$\mathbf a(t_0)\ne0$ implies $\mathbf a(t)\ne0$.
\end{proof}

For a single profile, this is the inviscid wave construction of
Craik and Criminale~\cite{CraikCriminale}. For the profiles considered
above, the curl and amplitude growth are
\[
 \curl\bigl(\mathbf aF(k\cdot x)\bigr)
       =(k\times\mathbf a)F'(k\cdot x),\qquad
 \tfrac12(|\mathbf a|^2)'=-\mathbf a\cdot D\mathbf a.
\]

\subsubsection{The preceding frame}\label{cc:sec:frame}
A wave evolving over a background matrix $B$ becomes part of the
background for the next wave. At the origin, its contribution to the
velocity gradient is $\mathbf a\otimes k$, and its unit wave vector and
velocity direction provide a moving orthonormal frame. We express the
older gradient $B$ in this frame to identify the couplings that drive the
next wave. The following calculation uses only $\tr B=0$; it requires no
symmetry of $B'+B^2$.

\begin{lemma}\label{cc:lem:frame}
Let $B(t)$ be a smooth trace-free matrix, and let $k\ne0$ and
$\mathbf a\ne0$ satisfy \eqref{cc:eq:wave-ode} with $D$ replaced by $B$.
Put $E_1=\frac{k}{|k|}$, $E_2=\frac{\mathbf a}{|\mathbf a|}$ and $E_3=E_1\times E_2$,
write $B_{ij}=E_i\cdot BE_j$ $(i,j=1,2,3)$ for the entries of $B$ in this
orthonormal frame, and let $\sigma=|k|\,|\mathbf a|$, $\nu=B_{12}$ and
$\eta=B_{32}$. Then
\begin{equation}\label{cc:eq:frame}
 E_1'=-\nu E_2-B_{13}E_3,\quad
 E_2'=\nu E_1-\eta E_3,\quad
 E_3'=B_{13}E_1+\eta E_2,
 \qquad(\log\sigma)'=B_{33}.
\end{equation}
\end{lemma}
\begin{proof}
Since $k'=-B^Tk$, we have $E_1'=-(\Id-E_1E_1^T)B^TE_1=-B_{12}E_2-B_{13}E_3$
and $(\log|k|)'=-B_{11}$. Since $\mathbf a'=-B\mathbf a+p_{\rm w} k$
with $p_{\rm w}|k|=2|\mathbf a|B_{12}$, we have
$E_2'=(\Id-E_2E_2^T)(-BE_2+2B_{12}E_1)=B_{12}E_1-B_{32}E_3$ and
$(\log|\mathbf a|)'=-B_{22}$. Then
$E_3'=E_1'\times E_2+E_1\times E_2'=B_{13}E_1+B_{32}E_2$, and
$(\log\sigma)'=-B_{11}-B_{22}=B_{33}$ because $\tr B=0$.
\end{proof}
